% !TEX root = ../main.tex
% !TEX program = lualatex

\documentclass[../main.tex]{subfiles}
\externaldocument{../main}

\begin{document}

\section{Continuous Process Improvement}

\subsection{Purpose}

\ac{cpi-process} is the disciplined use of improvement
methods to remove waste, reduce variation, elevate constraints, and improve
fleet outcomes.  Coursebook 3.7.3 links \ac{cpi-process} to the EDO tool set: Lean, Theory
of Constraints, and Six Sigma are practical methods for improving engineering,
maintenance, production, and administrative processes
\autocite{edo-3-7-3-intro-cpi-2025}.

\subsection{Lean thinking}

\begin{longtblr}[
  caption = {Five Lean principles.},
  label = {tab:lean-principles},
]{
  width = \linewidth,
  colspec = {X[0.75,l] X[1.7,l]},
  rowhead = 1,
  hlines,
  vlines,
}
\textbf{Principle} & \textbf{Meaning for an EDO process} \\
Value & Define value from the customer or fleet perspective. \\
Value stream & Map every step that creates, supports, delays, or obscures value. \\
Flow & Remove queues, rework, handoff delays, and avoidable batching. \\
Pull & Let demand signal the next step instead of pushing excess work into the system. \\
Perfection & Continue improving after the first visible bottleneck is removed. \\
\end{longtblr}

Lean waste is any activity that does not add value for the customer.  Common
waste categories include waiting, excess inventory, unnecessary motion,
transportation, overproduction, overprocessing, defects, and unused talent
\autocite{edo-3-7-3-intro-cpi-2025}.

\subsection{Theory of Constraints}

Theory of Constraints starts from the premise that every system has a limiting
constraint.  The five focusing steps are: identify the constraint, exploit the
constraint, subordinate other work to the constraint, elevate the constraint,
and repeat when the constraint moves
\autocite{edo-3-7-3-intro-cpi-2025}.  A bottleneck is a visible queue or
capacity limit; a constraint is the system limit that controls total throughput.
They may be the same, but the EDO should verify rather than assume.

\subsection{Six Sigma and DMAIC}

Six Sigma focuses on reducing variation and defects.  The common performance
shorthand is 3.4 defects per million opportunities.  The DMAIC improvement
cycle is define, measure, analyze, improve, and control
\autocite{edo-3-7-3-intro-cpi-2025}.

\begin{longtblr}[
  caption = {DMAIC board cue.},
  label = {tab:dmaic-board-cue},
]{
  width = \linewidth,
  colspec = {X[0.55,l] X[1.75,l]},
  rowhead = 1,
  hlines,
  vlines,
}
\textbf{Step} & \textbf{Question} \\
Define & What problem, customer, process, and outcome are in scope? \\
Measure & What data prove the current baseline and variation? \\
Analyze & What root causes explain the performance gap? \\
Improve & What countermeasures change the process? \\
Control & What controls prevent regression after the event ends? \\
\end{longtblr}

\paragraph{Board cue.}
Do not answer CPI as a slogan.  Name the method, state the constraint or waste,
show the metric, and explain how the control plan prevents backsliding.

\IfSubfilesClassLoaded{
  \RenewDocumentCommand{\entryname}{}{\textbf{\color{Modern} Acronym}}
  \RenewDocumentCommand{\descriptionname}{}{\textbf{\color{Modern} Definition}}
  \printnoidxglossary[
    type=\acronymtype,
    title=Acronyms,
    style=long-booktabs
  ]}{}
\end{document}
